\documentclass{article}
\usepackage{amsmath,amssymb}
\usepackage{xurl}
\usepackage[hidelinks]{hyperref}
% Shared commands for notes in docs/paper-gaps/.

\DeclareUnicodeCharacter{2080}{\ifmmode{}_{0}\else\textsubscript{0}\fi}
\DeclareUnicodeCharacter{2081}{\ifmmode{}_{1}\else\textsubscript{1}\fi}
\DeclareUnicodeCharacter{2082}{\ifmmode{}_{2}\else\textsubscript{2}\fi}
\DeclareUnicodeCharacter{2099}{\ifmmode{}_{n}\else\textsubscript{n}\fi}

\newcommand{\Lean}{\textsc{Lean}}
\ExplSyntaxOn
\NewDocumentCommand{\leanid}{m}{
  \ifmmode
    \textsf{\detokenize{#1}}
  \else
    \group_begin:
    \urlstyle{sf}
    \tl_set:Nn \l_tmpa_tl {#1}
    \tl_replace_all:Nnn \l_tmpa_tl {\_} {_}
    \exp_args:NV \path \l_tmpa_tl
    \group_end:
  \fi
}
\ExplSyntaxOff
\providecommand{\tr}{\operatorname{tr}}
\newcommand{\tnleanIssueBase}{https://github.com/LionSR/TNLean/issues/}
\newcommand{\tnleanPrBase}{https://github.com/LionSR/TNLean/pull/}
\newcommand{\tnleanDocBase}{https://sirui-lu.com/TNLean/docs/}
\newcommand{\ghissue}[1]{\href{\tnleanIssueBase#1}{GitHub issue \##1}}
\newcommand{\ghpr}[1]{\href{\tnleanPrBase#1}{PR \##1}}
\newcommand{\doclink}[1]{\href{\tnleanDocBase#1.html}{\path{#1}}}

% Dirac notation and the identity operator, as in blueprint/src/macros/common.tex.
% \providecommand keeps note-local definitions authoritative.
\usepackage{dsfont}
\providecommand{\ket}[1]{|#1\rangle}
\providecommand{\bra}[1]{\langle #1|}
\providecommand{\braket}[2]{\langle #1 | #2 \rangle}
\providecommand{\ketbra}[2]{|#1\rangle\!\langle #2|}
\providecommand{\Id}{\mathds{1}}
\providecommand{\one}{\mathds{1}}
\providecommand{\C}{\mathbb{C}}
\providecommand{\MN}[1]{M_{#1}(\C)}
\providecommand{\MD}{M_D(\C)}

% Verdict marker: \gapnote{<kind>}{<status>}. Kind is the policy
% classification (clarification, local-correction, scope-restriction,
% unfaithful, false-source, open-gap); status is open, wip (elimination
% underway), resolved, or historical. Machine-read by the site index;
% prints a verdict line.
\newcommand{\gapnote}[2]{\par\noindent\textbf{Verdict:} #1 (#2).\par}

\title{Weak-Hopf Averaging of the Mixed Path on its Active Physical Space}
\author{}
\date{2026-11-16}
\begin{document}
\maketitle
\gapnote{scope-restriction}{open}

\section{The source assertion}
Garre-Rubio, Lootens and Moln\'ar symmetrize parent interactions using a
canonical weak-Hopf integral in Section 5, lines 1276--1320, of
arXiv:2203.12563v3 \cite[Section~5]{GarreRubio2023MPOSym}.
Appendix B, lines 2463--2466, averages the support projection $P$, with
$E(P)=P$ and $E(I)=Q=\rho_2(1)$. Thus the literal average of the parent
is $Q-P$, whereas its ambient completion is $I-P$.
The distinction and the Bell-pair example with $Q=P$ are explained in
\href{https://sirui-lu.com/TNLean/paper-gaps/glm23_wha_parent_completion.pdf}
{Weak-Hopf Parent Terms: Support Averaging and Ambient Completion}.
The mixed interpolation of Section 5, lines 1584--1692, requires applying
this construction to its common symmetry.

\section{A unital realization cannot equal the enlarged tensor}
Write $a=(D_0+D_1)^2$ and let
$J:\mathbb C^a\longrightarrow\mathbb C^{a+d_0+d_1}$ include the first
summand. The common enlarged operator tensor is
$\widetilde T^{ij}_{\beta\alpha}=(J T_{\beta\alpha}J^*)_{ij}$.
Every physical row in either unused summand vanishes. In particular,
$\widetilde T^{ii}=0$ for such an $i$.

Suppose instead that $\widetilde T^{ij}=\psi(f_{ij})$, where
$f_{ij}(x)=\phi(x)_{ij}$, $\phi$ is a unital physical representation,
and $\psi$ is faithful on the convolution dual. Faithfulness would imply
$f_{ii}=0$, but $f_{ii}(1)=\phi(1)_{ii}=1$. This is a contradiction whenever
$d_0+d_1>0$. It uses neither fusion data nor any property of the integral.
It also does not concern adjoint closure: extension by zero can preserve
adjoint closure while destroying physical unitality.\footnote{
    The obstruction is stated in
    \leanid{MPSTensor.MPOSymmetry.arbitraryPhysicalMixedMPO_ne_weakHopfTensor},
    in \doclink{TNLean/MPS/Symmetry/MPOSymmetry/MixedPath/MPO/MixedWeakHopfRealizationObstruction}.
    The literal enlarged realization requested by \ghissue{10762} must therefore
    be changed. The independent adjoint-closure question \ghissue{10763} is not
    settled here.}

The conditional enlarged-space theorems in
\doclink{TNLean/MPS/Symmetry/MPOSymmetry/MixedPath/MPO/ArbitraryPhysicalMixedWeakHopf}
assume precisely the equality excluded above. For their parent-commutation
statements, injectivity and nonzero endpoint bond dimensions imply
$d_0,d_1>0$. Hence these hypotheses are incompatible: the conditional
implications are vacuous, rather than their conclusions contradicting a
proved result. No realization witness is constructed there. The weaker
boundary-closure and support-average statements are excluded by this
obstruction whenever either physical dimension is positive; the case
$d_0=d_1=0$ is not decided by this obstruction alone.

\section{The conditional construction on the active space}
Assume the two endpoint tensors are injective, their bond dimensions are
positive, and their actions have exact biorthogonal decompositions into
the same number of copies. Compress the endpoints to their polar support
coordinates and assemble the active mixed operator $T$.
Assume, additionally, that this actual $T$ has a weak-Hopf realization:
a complex coalgebra with multiplicative, star-preserving coproduct, a
unital physical star representation, a faithful convolution-dual
representation, and integral data normalized by
\begin{align}
    \sum\Delta(\Omega_{(1)})\Delta(T_\Omega(\Omega_{(2)})) & =\Delta(1).
    \label{eq:mixed_wha_normalization}
\end{align}
Here $T_\Omega$ denotes the specified linear pulling-through map; full
antipode axioms are not needed. No independent adjoint-closure hypothesis
is imposed. Neither existence of these data nor their construction from
F-symbols is asserted.

Let $h_\gamma$ be the continuous extended two-site interaction on the
active space, including its limiting values at $0$ and $1$. In the
interior it is $I-P_\gamma$ for the injective mixed tensor. Compatibility
and the normalized average give
\begin{align}
    E(h_\gamma)                          & =Q-I+h_\gamma, &
    \widehat h_\gamma:=E(h_\gamma)+(I-Q) & =h_\gamma.
    \label{eq:mixed_wha_completion}
\end{align}
The average is a finite sum of fixed left and right multiplications,
so it is continuous. Both identities in~(\ref{eq:mixed_wha_completion})
therefore extend to $[0,1]$ by density of $(0,1)$.
If the weak-Hopf algebra is finite dimensional, every boundary operator
is represented by an algebra element. Hence $Q$ commutes with it, and
the literal active average also commutes with every boundary operator. This argument does not
identify the limiting support at either endpoint with that endpoint
tensor's canonical support.

For $J_2=J\otimes J$, the completed enlarged interaction is
\begin{align}
    \widetilde h_\gamma
     & =I-J_2(I-\widehat h_\gamma)J_2^*
    =I-J_2J_2^*+J_2\widehat h_\gamma J_2^*.
    \label{eq:mixed_wha_enlargement}
\end{align}
By~(\ref{eq:mixed_wha_completion}), this is exactly the existing enlarged
mixed interaction, hence positive, continuous, and uniformly gapped on
periodic rings of all lengths $N\ge2$. The same equality preserves its
already established ground spaces. Only the completed interaction is
covered by these claims; the literal uncompleted average can have extra
zero-energy states.

Adjoint closure on the active space follows from the realization.
The same adjoint boundary works after inclusion, since
$O_{\widetilde T,X}^N=J_N O_{T,X}^N J_N^*$.
Compatibility then gives local commutation in the interior; continuity
extends it to the endpoints. Local commutation implies commutation of the
periodic sum with the periodic common MPO at every $N\ge2$.\footnote{
    The conditional average identities and active-space commutation are in
    \doclink{TNLean/MPS/Symmetry/MPOSymmetry/MixedPath/MPO/MixedWeakHopfAverage}.
    The enlarged construction, positivity, continuity, periodic commutation and
    uniform gap are in
    \doclink{TNLean/MPS/Symmetry/MPOSymmetry/MixedPath/MPO/ArbitraryPhysicalMixedWeakHopfParent}.
    The gap is transferred by exact equality with the previously proved path,
    not by claiming that averaging preserves a spectral gap.}

\section{Remaining reconstruction problem}
The enlarged-space unital realization is impossible as specified, not
merely unproved. A corrected program must construct the normalized
realization on the active space from suitably admissible fusion and
star data, or formulate and justify a nonunital physical representation
on the enlarged space. The former realization is still an additional
hypothesis here. General exact action and fusion decompositions alone
do not supply physical star data. This note does not assert a
counterexample to the source's phase-classification theorem, nor close
the reconstruction issue.
\bibliographystyle{alpha}
\bibliography{references}
\end{document}
